% Source: MQ33b_v1.html
\documentclass[11pt]{article}
\usepackage[utf8]{inputenc}
\usepackage[T1]{fontenc}
\usepackage[margin=1in]{geometry}
\usepackage{amsmath}
\usepackage{amssymb}
\usepackage{siunitx}
\usepackage{enumitem}
\usepackage{parskip}

\title{PS160 -- MQ33b\_v1.html}
\date{}

\begin{document}
\maketitle

\section*{Questions}

\begin{enumerate}[leftmargin=*]
  \item The table below shows the index of refraction for BK7 optical glass at several different wavelengths. White light is incident on the top face of a large block of this glass at an angle of 49$^{\circ}$ with respect to the normal. If the block is 0.3 m thick, by how many millimeters are beams at wavelengths 2 and 4 separated at the bottom face of the block?

Number
Wavelength (nm)
Index of Refraction

1
400
1.531

2
475
1.523

3
550
1.519

4
625
1.515 \hfill \textit{(1 pts, Numeric)}

  \item A ray of light is traveling in a fluid having an index of refraction 2 when it encounters a boundary with another fluid that has index of refraction 1.4. Calculate the critical angle in degrees for total internal reflection. \\ \hfill \textit{(1 pts, Numeric)}

  \item A liquid-air boundary has a critical angle of 44.1$^{\circ}$ for total internal reflection. If a ray of light traveling in air has an incident angle of 28$^{\circ}$ at the liquid-air boundary, calculate the angle between the refracted ray and the surface normal. \hfill \textit{(1 pts, Numeric)}

  \item Monochromatic, polarized light passes through a polarization filter with a polarization axis that is rotated 16$^{\circ}$ with respect to the light polarization. The light then passes through a second polarization filter with axis rotated a further 47$^{\circ}$ from that of the previous filter. Calculate the ratio of the final light intensity to the original light intensity. \hfill \textit{(1 pts, Numeric)}

  \item A ray of light traveling through air reflects off the surface of an unknown fluid. When the ray's angle of incidence is 47.2$^{\circ}$, the light is completely polarized. Calculate the fluid's index of refraction. \hfill \textit{(1 pts, Numeric)}

\end{enumerate}

\end{document}